\documentclass[allauthors,dutch]{../../../cursuspresentatie3}


\setbeamertemplate{headline}
{%
    \placetarget
    \let\placetarget\relax
    \begin{beamercolorbox}[colsep=1.5pt]{upper separation line head}%
    \end{beamercolorbox}
    \begin{beamercolorbox}{section in head/foot}
        \vskip2pt
        \usebeamerfont{section in head/foot}\usebeamercolor[fg]{section in head/foot}%
        \color{fg!80!bg}
        \hfil
        \adjustbox{
            set depth=2pt,
            padding={5pt 5pt 5pt 5pt},
            rndframe={color=navigationTileOutline,width=1pt}{3pt 3pt 3pt 3pt},margin={10pt 0pt 0pt 0pt}
        }{
            \hyperlink{herhaling}{\customhead{1}{Korte herhaling}}
        }%
        \hfil
        \adjustbox{
            set depth=2pt,
            padding={5pt 5pt 5pt 5pt},
            rndframe={color=navigationTileOutline,width=1pt}{3pt 3pt 3pt 3pt},margin={10pt 0pt 0pt 0pt}
        }{%
            \hyperlink{newcommand}{\customhead{2}{Newcommand}}
        }
        \hfil
        \adjustbox{
            set depth=2pt,
            padding={5pt 5pt 5pt 5pt},
            rndframe={color=navigationTileOutline,width=1pt}{3pt 3pt 3pt 3pt},margin={10pt 0pt 0pt 0pt}
            % rndcorners={3pt 0pt 3pt 0pt}
        }{%
            \hyperlink{comments}{\customhead{3}{Comments}}
        }%
        \hfil
        \adjustbox{
            set depth=2pt,
            padding={5pt 5pt 5pt 5pt},
            rndframe={color=navigationTileOutline,width=1pt}{3pt 3pt 3pt 3pt},margin={10pt 0pt 0pt 0pt}
            % rndcorners={3pt 0pt 3pt 0pt}
        }{%
            \hyperlink{theorem_environment}{\customhead{4}{Theorem environment}}
        }%
        \hfil
        \adjustbox{
            set depth=2pt,
            padding={5pt 5pt 5pt 5pt},
            rndframe={color=navigationTileOutline,width=1pt}{3pt 3pt 3pt 3pt},margin={10pt 0pt 0pt 0pt}
            % rndcorners={3pt 0pt 3pt 0pt}
        }{%
            \hyperlink{packages}{\customhead{5}{Packages}}
        }%
        \hfil  
        \adjustbox{
            set depth=2pt,
            padding={5pt 5pt 5pt 5pt},
            rndframe={color=navigationTileOutline,width=1pt}{3pt 3pt 3pt 3pt},margin={10pt 0pt 0pt 0pt}
        }{%

            \hyperlink{biblio}{\customhead{6}{Bibliografie}}%
        }%
        \hfil
        \vskip4pt
    \end{beamercolorbox}%
    \begin{beamercolorbox}[colsep=1.5pt]{lower separation line head}%
    \end{beamercolorbox}%
}

\usefonttheme{professionalfonts} % reset math mode font to default
\def\importslidei#1#2{%
	\import{../../../slides/ivo/#1}{#2}
}

\def\importslideh#1#2{%
	\import{../../../slides/Hylke/#1}{#2}
}

\def\importslidej#1#2{%
	\import{../../../slides/Jagna/2025-2026/#1}{#2}
}

\def\importslidey#1#2{%
	\import{../../../slides/yoram/#1}{#2}
}

\newcommand{\cmfontblock}[1]{%
  {\rmfamily\selectfont #1}%
}

\title{\LaTeX{}-cursus 2}
\author{\TeX niCie}
\date{6 oktober 2026}

% Als je het bestand dumpt naar een format, worden packages hierna niet meegedumpt maar elke keer
% fris ingeladen
\csname endofdump\endcsname

\usepackage{xcolor}
\setminted[tex]{fontsize=\small, autogobble=true, linenos=true, frame=single, escapeinside=||}
\setminted[bib]{fontsize=\small, autogobble=true, linenos=false, frame=none, escapeinside=||}
\usepackage{tcolorbox}
% \usepackage{}
\definecolor{darkgreen}{rgb}{0.2,0.7,0.11}
%\makeatletter
%\patchcmd{\@thm}{\thm@headfont{\scshape}}{\thm@headfont{\scshape\bfseries}}{}{}
%\patchcmd{\@thm}{\thm@notefont{\fontseries\mddefault\upshape}}{}{}{}
%\makeatother
%\newtheorem{mylemma}[]{Lemma}

\begin{document}

\importslideh{beginners_NL}{welcome.tex}

\section{Korte herhaling}
\def\placetarget{\hypertarget{herhaling}{}}

\importslidei{2025-eerstejaars_cursus_1}{Math_mode_display.tex}

% slides van Yoram
\section{Newcommand}
\def\placetarget{\hypertarget{newcommand}{}}

\importslidey{2025-eerstejaars_cursus_2}{newcommand00.tex}
\importslidey{2025-eerstejaars_cursus_2}{newcommand01.tex}
\importslidey{2025-eerstejaars_cursus_2}{newcommand02.tex}
\importslidey{2025-eerstejaars_cursus_2}{newcommand03.tex}
\importslidey{2025-eerstejaars_cursus_2}{newcommand04.tex}
% \importslidey{2025-eerstejaars_cursus_2}{newcommand05.tex}
\importslidey{2025-eerstejaars_cursus_2}{newcommand06.tex}

\section{Comments}
\def\placetarget{\hypertarget{comments}{}}

\importslidey{2025-eerstejaars_cursus_2}{comments00.tex}
\importslidey{2025-eerstejaars_cursus_2}{comments01.tex}
\importslidey{2025-eerstejaars_cursus_2}{comments02.tex}

\section{Theorem environment}
\def\placetarget{\hypertarget{theorem_environment}{}}
% Was niet nodig voor natuurkunde?
\importslidey{2025-eerstejaars_cursus_2}{theorems00.tex}
\importslidey{2025-eerstejaars_cursus_2}{theorems01.tex}
\importslidey{2025-eerstejaars_cursus_2}{theorems02.tex}
\importslidey{2025-eerstejaars_cursus_2}{theorems03.tex}
\importslidey{2025-eerstejaars_cursus_2}{theorems04.tex}
\importslidey{2025-eerstejaars_cursus_2}{theorems05.tex}

\section{Packages}
\def\placetarget{\hypertarget{packages}{}}
\importslidej{cursus2}{2025_eerstejaarscursus2_packages.tex}

\section{Bibliografie}
\def\placetarget{\hypertarget{biblio}{}}
\importslideh{bib}{bib-introduction.tex}
\importslideh{bib}{bib-structure.tex} % 6
\importslideh{bib}{bib-bibentry.tex} %13
\importslideh{bib}{bib-bibfilestelen.tex}
\importslideh{bib}{bib-mult_authors.tex}

%\importslideh{bib}{bib-cite.tex}
\importslidei{2025-scriptie_cursus}{Hyperref_cleveref.tex}

\importslidej{cursus2}{2025_eerstejaarscursus2_zoekopdrachten.tex}


\section{Afsluiting}
\begin{frame}
	\frametitle{Afsluiting}
	Leuk dat jullie er waren!
	
	Voor vragen zijn wij bereikbaar op 
	\href{mailto:texnicie@a-eskwadraat.nl}{texnicie@a-eskwadraat.nl}
\end{frame}

\end{document}